\section{Ablaciones intra-sujeto de preprocesamiento y aumentación}
\label{sec:ablaciones-f4}

Ningún efecto de preprocesamiento o aumentación es común a las cuatro arquitecturas: la ventana corta, la
banda $\mu+\beta$ y la reflexión de canales se asocian con una exactitud mayor en algunas y no en otras
(tabla~\ref{tab:ablaciones-f4}). Los resultados de \eegnet{} y el Conformer están \emph{condicionados al
régimen de entrenamiento}, porque no existe una base con recorte con la que estimar una cota
(sección~\ref{sec:confusor}); en LaBraM la cota es de $0.0182$, y el contraste de EEGMamba no cambia de
recorte, aunque el desplazamiento de $+0.009$ de esa arquitectura entre semillas es del orden de varios
de los efectos de la tabla. Todas las variantes tienen una sola semilla, por lo que las equivalencias valen
para la semilla $2026$ (el desplazamiento medio entre semillas llega a $0.008$ y $0.009$). Los efectos se
describen como asociaciones: no se contrastó la interacción entre arquitectura y variante.

\begin{table}[H]
\centering
\caption{Ablaciones intra-sujeto: diferencia de exactitud de la variante menos la base de la semilla
$2026$, sobre los $54$ participantes, con intervalo al $95$\,\% de percentiles y $p_{\text{Holm}}$ dentro de la
familia de la arquitectura (ocho contrastes en \eegnet{} y Conformer, tres en LaBraM y EEGMamba). Margen
mín.: menor margen con el que se cumpliría la equivalencia (equivalente dentro de $\pm 0.03$ si
es $\leq 0.030$). Lectura: \emph{sin cambio de recorte} (EEGMamba); \emph{sin cota
estimable}, condicionado al régimen (\eegnet{} y Conformer); \emph{supera la cota} o \emph{no
interpretable}, según la magnitud frente a la cota de $0.0182$ (LaBraM).}
\label{tab:ablaciones-f4}
\footnotesize
\setlength{\tabcolsep}{3pt}
\resizebox{\textwidth}{!}{%
\begin{tabular}{llccccl}
\toprule
\textbf{Arq.} & \textbf{Variante} & Dif. media & IC $95$\,\% & $p_{\text{Holm}}$ & Margen mín. & Lectura \\
\midrule
\eegnet{} & banda $\mu$ ($8$ a $13$\,Hz) & $-0.0300$ & $[-0.0635,\ +0.0032]$ & 0.49 & 0.058 & sin cota estimable \\
 & banda $\beta$ ($13$ a $30$\,Hz) & $-0.0456$ & $[-0.0788,\ -0.0098]$ & 0.009 & 0.073 & sin cota estimable \\
 & banda $\mu+\beta$ ($8$ a $30$\,Hz) & $-0.0060$ & $[-0.0450,\ +0.0342]$ & 0.90 & 0.039 & sin cota estimable \\
 & ventana $0$ a $2$\,s & $+0.0417$ & $[+0.0294,\ +0.0541]$ & $<0.001$ & 0.052 & sin cota estimable \\
 & ventana $0$ a $3$\,s & $+0.0045$ & $[-0.0054,\ +0.0145]$ & 0.83 & 0.013 & sin cota estimable \\
 & sin aumentación & $+0.0073$ & $[-0.0046,\ +0.0200]$ & 0.79 & 0.018 & sin cota estimable \\
 & reflexión de canales & $+0.0260$ & $[+0.0120,\ +0.0420]$ & 0.005 & 0.039 & sin cota estimable \\
 & ruido en Fourier & $+0.0019$ & $[-0.0035,\ +0.0074]$ & 0.90 & 0.007 & sin cota estimable \\
\midrule
Conformer & banda $\mu$ ($8$ a $13$\,Hz) & $+0.0459$ & $[+0.0084,\ +0.0850]$ & 0.23 & 0.078 & sin cota estimable \\
 & banda $\beta$ ($13$ a $30$\,Hz) & $+0.0355$ & $[+0.0008,\ +0.0723]$ & 0.45 & 0.066 & sin cota estimable \\
 & banda $\mu+\beta$ ($8$ a $30$\,Hz) & $+0.0781$ & $[+0.0407,\ +0.1165]$ & 0.002 & 0.111 & sin cota estimable \\
 & ventana $0$ a $2$\,s & $-0.0044$ & $[-0.0193,\ +0.0090]$ & 1.00 & 0.017 & sin cota estimable \\
 & ventana $0$ a $3$\,s & $-0.0097$ & $[-0.0175,\ -0.0022]$ & 0.13 & 0.016 & sin cota estimable \\
 & sin aumentación & $-0.0068$ & $[-0.0126,\ -0.0009]$ & 0.23 & 0.012 & sin cota estimable \\
 & reflexión de canales & $+0.0219$ & $[+0.0120,\ +0.0318]$ & 0.002 & 0.030 & sin cota estimable \\
 & ruido en Fourier & $+0.0010$ & $[-0.0049,\ +0.0071]$ & 1.00 & 0.006 & sin cota estimable \\
\midrule
LaBraM & ventana $0$ a $2$\,s & $-0.0217$ & $[-0.0406,\ -0.0030]$ & 0.12 & 0.038 & supera la cota \\
 & reflexión de canales & $-0.0139$ & $[-0.0265,\ -0.0011]$ & 0.043 & 0.024 & no interpretable \\
 & ruido en Fourier & $-0.0031$ & $[-0.0098,\ +0.0037]$ & 0.25 & 0.009 & no interpretable \\
\midrule
EEGMamba & ventana $0$ a $2$\,s & $+0.0242$ & $[+0.0089,\ +0.0390]$ & 0.004 & 0.037 & sin cambio de recorte \\
 & reflexión de canales & $+0.0232$ & $[+0.0090,\ +0.0381]$ & 0.015 & 0.036 & sin cambio de recorte \\
 & ruido en Fourier & $-0.0069$ & $[-0.0150,\ +0.0012]$ & 0.12 & 0.014 & sin cambio de recorte \\
\bottomrule
\end{tabular}
}
\end{table}

\paragraph{Ventana.} La ventana de \seg{0} a \seg{2} (la de \seg{0} a \seg{4} se fijó sin evaluar
alternativas) da una exactitud mayor en \eegnet{} ($+0.0417$; $46$ de $54$ participantes mejoran,
$r = 0.68$, $p_{\text{Holm}} < 0.001$) y en EEGMamba ($+0.0242$, $p_{\text{Holm}} = 0.004$), es
equivalente a la base en el Conformer ($-0.0044$) y en LaBraM se asocia con una exactitud menor
($-0.0217$, sobre la cota de $0.0182$, $p_{\text{Holm}} = 0.12$), sin ser concluyente. La ventana de
\seg{0} a \seg{3} es equivalente a la de \seg{0} a \seg{4} en \eegnet{} y el Conformer. La ventana del
resto del documento sigue siendo la fijada de antemano.

\paragraph{Subbandas.} La banda $\mu+\beta$ se asocia con una exactitud mayor en el Conformer
($+0.0781$; $39$ de $54$ participantes mejoran, $r = 0.50$, $p_{\text{Holm}} = 0.002$), y la banda $\beta$
con una menor en \eegnet{} ($-0.0456$, $p_{\text{Holm}} = 0.009$); las demás no sobreviven a la
corrección. El signo es opuesto en las dos arquitecturas evaluadas y no se confirmó con otra semilla; no se
evaluaron en LaBraM ni en EEGMamba.

\paragraph{Aumentación.} La reflexión de canales se asocia con una exactitud mayor en \eegnet{}
($+0.0260$, $p_{\text{Holm}} = 0.005$), el Conformer ($+0.0219$, $0.002$) y EEGMamba ($+0.0232$, $0.015$),
y no es interpretable en LaBraM ($-0.0139$, bajo la cota de $0.0182$; $-0.0118$ con la base del mismo
régimen). El ruido en Fourier es equivalente a la base en las cuatro (margen mínimo de $0.006$ a $0.014$),
con equivalencia dependiente del régimen en las tres confundidas. Quitar toda la aumentación cambia la
exactitud en menos de un punto en \eegnet{} ($+0.0073$) y el Conformer ($-0.0068$): la aumentación por
defecto no aporta ni resta de forma detectable. Con la familia secundaria de $22$ contrastes dejan de ser
significativas la reflexión de canales en EEGMamba ($0.118$) y en LaBraM ($0.217$).

\paragraph{La parada temprana como explicación alternativa.} En $21$ de los $22$ contrastes el signo de la
diferencia de exactitud coincide con el del cambio en el porcentaje de particiones cuyo mejor modelo es el
de la época $30$, el suelo del esquema (la exactitud sube cuando el porcentaje baja;
tabla~\ref{tab:parada-f4}). Por ejemplo, la ventana de \seg{0} a \seg{2} reduce ese porcentaje en \eegnet{}
del $16.6$ al $7.6$\,\%. Una parte de los efectos puede reflejar si el entrenamiento progresa más allá del
calentamiento con $144$ ensayos y no la calidad de la representación; el diseño no separa ambas lecturas,
aunque en EEGMamba el $79.4$\,\% de las particiones con ventana corta se sigue deteniendo en el suelo, de
modo que la parada temprana no explica por sí sola esa mejora. Este análisis es exploratorio y no
pertenece a ninguna familia de contrastes.

\begin{table}[H]
\centering
\caption{Porcentaje de las $1\,080$ particiones intra-sujeto cuyo mejor modelo es el de la época $30$ (el
suelo: diez épocas de calentamiento y veinte de paciencia), por arquitectura y variante, semilla $2026$.}
\label{tab:parada-f4}
\small
\begin{tabular}{lcccc}
\toprule
 & \textbf{\eegnet{}} & \textbf{Conformer} & \textbf{LaBraM} & \textbf{EEGMamba} \\
\midrule
Base (semilla 2026) & 16.6\,\% & 53.1\,\% & 21.1\,\% & 85.6\,\% \\
Ventana $0$ a $2$\,s & 7.6\,\% & 63.1\,\% & 28.5\,\% & 79.4\,\% \\
Reflexión de canales & 12.2\,\% & 28.1\,\% & 26.9\,\% & 58.4\,\% \\
Ruido en Fourier & 15.2\,\% & 55.3\,\% & 25.4\,\% & 87.2\,\% \\
Ventana $0$ a $3$\,s & 13.8\,\% & 61.9\,\% & -- & -- \\
Sin aumentación & 12.0\,\% & 58.9\,\% & -- & -- \\
Banda $\mu$ & 22.4\,\% & 29.2\,\% & -- & -- \\
Banda $\beta$ & 31.4\,\% & 28.4\,\% & -- & -- \\
Banda $\mu+\beta$ & 22.8\,\% & 21.9\,\% & -- & -- \\
\bottomrule
\end{tabular}
\end{table}

